Bolivia opera con dos sistemas de precios del dólar: el tipo de cambio
oficial administrado por el BCB y el mercado paralelo (\emph{dólar blue})
que se forma en plataformas P2P de criptomonedas. Este informe documenta,
siguiendo la metodología CRISP-DM \cite{chapman2000crisp},
las fases de comprensión del negocio, comprensión de los datos, preparación
y modelado, cuyo objetivo es \textbf{predecir} el nivel del paralelo y la
brecha oficial--paralelo a horizontes de una hora y un día. Se describe el pipeline
que recolecta y normaliza las cotizaciones de la API pública de
dolarbluebolivia.click y de la Gaceta Oficial de Bolivia, el análisis
exploratorio de la brecha, la tabla de características y los modelos de
previsión (regresión lineal, OLS con errores HAC y ARIMA) evaluados contra
una línea base naïve en una ventana de prueba de las últimas 168~h (una semana), con
frecuencia horaria.

\medskip
\noindent\textit{Datos al: 2026-10-24 02:33 UTC
\unskip.}
